% =============================================================================
% SECCIÓN 4: DISCUSIÓN
% Capítulo 2: Instrumentación IoT y Caracterización Metrológica
% =============================================================================

% -----------------------------------------------------------------------------
% 4.1. SÍNTESIS Y POSICIONAMIENTO METROLÓGICO
% -----------------------------------------------------------------------------
Los resultados obtenidos validan que un nodo de sensores NDIR de bajo
costo, sometido a protocolos de calibración de dos puntos y compensación
cruzada por temperatura y humedad, alcanza un desempeño metrológico
aceptable para la cuantificación de \ce{CO2} y \ce{CH4} en biorreactores
de \textit{H. illucens} bajo condiciones de alta humedad relativa. El
MAPE de 8.3\,\% para \ce{CO2} y 11.7\,\% para \ce{CH4} cumple con el
umbral de aceptación establecido en la hipótesis H1 (MAPE $\leq 15$\,\%)
y sitúa el sistema dentro del rango de precisión requerido para
aplicaciones de monitoreo ambiental operativo, aunque por debajo de la
exactitud de analizadores de referencia de laboratorio (MAPE $< 2$\,\%
para GC-TCD).

% -----------------------------------------------------------------------------
% 4.2. COMPARACIÓN CON ESTRATEGIAS DE CALIBRACIÓN EN LA LITERATURA
% -----------------------------------------------------------------------------
La literatura reciente ha abordado la calibración de sensores de bajo
costo mediante técnicas de aprendizaje automático. \textcite{biagiDevelopmentMachineLearningbased2024} desarrollaron estaciones
multiparamétricas para \ce{CO2} y \ce{CH4} en aire ambiente, utilizando
redes neuronales para corrección post-hoc de la deriva térmica y
hídrica. Si bien este enfoque reduce el error residual en entornos
controlados, introduce una dependencia computacional que dificulta su
implementación en dispositivos edge de recursos limitados. En contraste,
el modelo de compensación física implementado en este capítulo
(Ecuación~\ref{eq:compensacion_ndir}) opera enteramente en el
microcontrolador ESP32, con un costo computacional de $< 5$\,ms por
ciclo de corrección y una memoria RAM requerida de $< 2$\,KB. Esta
eficiencia es crítica para contextos agroindustriales del Sur Global,
donde la conectividad intermitente y la energía restringida hacen
inviable el procesamiento en la nube \autocite{duobieneDevelopmentWirelessSensor2022}.

Por otro lado, trabajos previos en monitoreo de sistemas de BSF han
priorizado variables ambientales de confort (temperatura, humedad, pH)
sin validación metrológica rigurosa de los gases de efecto invernadero.
\textcite{vanIntegrationInternetofThingsSustainable2022} implementaron
un sistema IoT para cría de BSF con sensores DHT22 y pH, pero no
integraron sensores NDIR validados ni reportaron incertidumbre
expandida. El aporte distintivo de este capítulo radica en la
cuantificación explícita de la incertidumbre expandida ($U$, $k=2$)
y la trazabilidad a patrones certificados, elementos indispensables
para que los datos generados sean admisibles en protocolos de
Medición, Reporte y Verificación (MRV) de emisiones.

% -----------------------------------------------------------------------------
% 4.3. MECANISMOS DE DERIVA EN CONDICIONES AGRESIVAS
% -----------------------------------------------------------------------------
La deriva observada en HR $> 80$\,\% se explica principalmente por dos
mecanismos físicos. En el sensor de \ce{CO2}, la condensación de vapor
sobre la óptica del detector reduce la transmitancia de la banda de
absorción a 4.26\,$\mu$m, generando una lectura positiva sistemática
que se acumula con el tiempo. La membrana ePTFE implementada mitigó
este efecto al actuar como barrera líquida-permeable a gases, pero su
eficacia disminuye cuando la humedad alcanza la saturación relativa
sobre la superficie de la membrana. En el sensor de \ce{CH4}, la deriva
negativa se atribuye a la interferencia espectral de compuestos
orgánicos volátiles (COV) emitidos durante la descomposición del
sustrato proteico (ácido butírico, mercaptanos), los cuales presentan
bandas de absorción superpuestas en la región de 3.3\,$\mu$m
\autocite{chenEffectMoistureContent2019}. Este fenómeno explica por qué
el MAPE de \ce{CH4} (11.7\,\%) supera al de \ce{CO2} (8.3\,\%) y
sugiere que, para aplicaciones que requieran cuantificación precisa de
metano a niveles $< 500$\,ppm, sería necesario incorporar un sensor NDIR
de doble longitud de onda o un filtro de corrección espectral por
interferentes.

% -----------------------------------------------------------------------------
% 4.4. INTERPRETACIÓN DE LAS DIFERENCIAS ENTRE DIETAS
% -----------------------------------------------------------------------------
La marcada diferencia en el perfil de \ce{CH4} entre la dieta D1
(proteica) y la D4 (fibrosa) tiene implicaciones directas para el
diseño de protocolos de monitoreo. La dieta D1, con mayor carga
biodegradable y relación C:N baja, favoreció la formación rápida de
núcleos anaerobios en el interior del sustrato denso, activando la
metanogénesis bacteriana tras un periodo de latencia de 3--8 minutos
durante el confinamiento hermético. Este comportamiento de umbral
valida empíricamente el modelo teórico de \textcite{bekkerImpactSubstrateMoisture2021a},
quienes reportaron que la humedad del sustrato $> 70$\,\% es un
detonante crítico de microzonas anaeróbicas en sistemas BSF. En
contraste, la dieta D4, de mayor porosidad y menor densidad energética,
permitió una aireación más eficiente del lecho, retrasando el inicio de
la producción de metano y reduciendo su concentración máxima en un
factor de 1.4.

Desde la perspectiva del monitoreo operativo, estos hallazgos indican
que el umbral de alerta para anaerobiosis no puede ser un valor fijo
de \ce{CH4}, sino que debe adaptarse al tipo de sustrato. Para dietas
de alta carga proteica (D1), una concentración de \ce{CH4} $> 1000$\,ppm
en los primeros 10 min de confinamiento constituye un indicador temprano
de saturación metabólica, mientras que para dietas fibrosas (D4) el
mismo umbral podría ser inadecuado o generar falsos negativos.

% -----------------------------------------------------------------------------
% 4.5. LIMITACIONES DEL ESTUDIO
% -----------------------------------------------------------------------------
Este capítulo presenta cuatro limitaciones principales que condicionan
la generalización de los resultados. Primera, el diseño experimental se
ejecutó a escala de mesocosmos de laboratorio (5.2\,L), con ciclos de
confinamiento hermético de duración limitada ($<< 30$\,min). Aunque
este enfoque permite el control metrológico riguroso de las
concentraciones, no replica la dinámica de un biorreactor industrial
abierto o de flujo continuo, donde la ventilación forzada modifica los
gradientes de gas en el espacio de cabeza. Segunda, el monitoreo no
fue continuo durante los 15 días del ciclo de bioconversión, sino
puntual en ventanas de medición. Esta limitación se aborda
sistemáticamente en los Capítulos \ref{sec:cap2_discusion} y 4,
donde el modelo mecanicista y el sistema híbrido interpolan y predicen
las emisiones entre mediciones. Tercera, el sensor de \ce{CH4} empleado
tiene un límite de detección efectivo de $\approx 200$\,ppm, lo cual
impide cuantificar concentraciones basales o eventos de muy baja
emisión que podrían ser relevantes para balances de carbono de alta
resolución. Cuarta, no se incluyó la medición directa de \ce{NH3} ni
\ce{N2O} por limitaciones de costo en sensores electroquímicos
selectivos, aunque la arquitectura modular del nodo permite su
integración futura sin rediseño estructural.

% -----------------------------------------------------------------------------
% 4.6. IMPLICACIONES PARA LA GESTIÓN AGROINDUSTRIAL EN EL SUR GLOBAL
% -----------------------------------------------------------------------------
El sistema validado en este capítulo tiene un potencial tangible de
transferencia a productores agroindustriales de pequeña y mediana escala
en el Valle del Cauca y regiones tropicales análogas. El costo marginal
del nodo IoT (estimado en $< 80$\,USD en componentes) representa menos
del 2\,\% del valor del producto potencial (biomasa larval + frass) en
un ciclo de 15 días, haciendo viable su adopción económica. La
detección temprana de picos de \ce{CH4} permite intervenir el proceso
mediante volteo del sustrato o ajuste de la densidad larval antes de
que la eficiencia de conversión se degrade irreversiblemente. Desde una
perspectiva de gobernanza ambiental, la disponibilidad de datos con
trazabilidad metrológica habilita la participación de estos productores
en mecanismos de economía circular con credenciales verificables,
cerrando la brecha entre la innovación tecnológica y el acceso a
mercados de carbono \autocite{jenkinsProcessingPoultryManure2025}.

% -----------------------------------------------------------------------------
% 4.7. CIERRE Y TRANSICIÓN
% -----------------------------------------------------------------------------
En síntesis, este capítulo establece que la instrumentación IoT de
bajo costo puede alcanzar estándares de calidad datos admisibles para
modelado predictivo, siempre que se incorporen protocolos de
calibración, compensación cruzada y validación contra métodos de
referencia. El dataset experimental generado ---con series temporales
de \ce{CO2} y \ce{CH4} validadas, caracterizadas por dieta y
condicionadas por variables ambientales--- constituye la base de
entrenamiento y validación para el modelo mecanicista DEB desarrollado
en el Capítulo \ref{chap:deb} y para la arquitectura híbrida HNODE
presentada en el Capítulo \ref{chap:hnode}. La precisión alcanzada
(MAPE $< 12$\,\%) garantiza que las señales de entrada a dichos modelos
no introduzcan un error instrumental que supere la incertidumbre
biológica inherente al sistema.
